% assets/a3-wedge.svg — 45 도 한계가 요구 공간에 만드는 쐐기와 null space 보정
%
%   bash _tools/tikz2svg.sh assets/src/a3-wedge.tex assets/a3-wedge.svg
%
% 이 그림이 말하는 것
%   추진기 하나의 각도 한계는 (F_x, F_y) 평면에서 **나비넥타이 꼴 영역**이다.
%       |delta| <= dmax   <=>   |F_y| <= tan(dmax) |F_x|
%   이것을 pseudo-inverse 해에 대입하면, 요구 (X, Y) 평면에 **쐐기 네 개**가 생긴다.
%   N = 0 단면에서 1.3111 < X/Y < 3.3111 이 안 되는 구간이고,
%   각도로는 원점에서 16.80 도 ~ 37.33 도 사이 방향이다.
%
%   쐐기는 pseudo-inverse **한 점**을 썼을 때의 이야기다. 같은 tau 를 내는 해가
%   직선 하나를 이루므로, 그 직선 위를 미끄러지면(null space 보정) 상당수가 되살아난다.
%   격자 실측 — 낼 수 있는 비율 19.7 % -> 25.3 %
%
% 좌표 주의
%   (나)는 **요구 공간**이다. 배를 위에서 본 그림이 아니다.
%   가로축 X (서지 힘), 세로축 Y (스웨이 힘). 방향만 보는 그림이라 눈금은 비율이다

\documentclass[border=6pt]{standalone}
% capstone-style.tex — 이 볼트의 TikZ 그림이 공유하는 스타일
%
% 쓰는 법:  % capstone-style.tex — 이 볼트의 TikZ 그림이 공유하는 스타일
%
% 쓰는 법:  % capstone-style.tex — 이 볼트의 TikZ 그림이 공유하는 스타일
%
% 쓰는 법:  \input{capstone-style}   (assets/src/ 안에서)
% 빌드:     bash _tools/tikz2svg.sh assets/src/w02-node-topic.tex assets/w02-node-topic.svg
%
% 왜 TikZ 인가
%   손으로 SVG 좌표를 쓰면 화살촉·선 굵기·글자 크기가 그림마다 미묘하게 달라진다.
%   그것이 "자료가 어설퍼 보인다"의 정체다. 스타일을 한 번만 정하고 상대좌표로
%   배치하면 좌표를 고쳐도 그림이 무너지지 않는다.
%   수식은 본문과 같은 TeX 글꼴로 나온다.
%
% 한글
%   ko.TeX(DVI 경로)를 쓴다. XeLaTeX 은 TikZ 그래픽을 PDF special 로 내보내는데
%   dvisvgm 이 그것을 받으려면 Ghostscript < 10.01 이 필요하다. 이 PC 에는
%   10.07 이 깔려 있어 그 경로가 막힌다. latex -> DVI -> dvisvgm 이 유일하게
%   동작하는 조합이다.

\usepackage[hangul]{kotex}
\usepackage{tikz}
\usepackage{amsmath}
\usepackage{xcolor}
\usetikzlibrary{arrows.meta, positioning, calc, fit, backgrounds, decorations.pathmorphing}

% ---- 색 --------------------------------------------------------------------
% 강조는 하나만. 나머지는 검정. 흑백 인쇄에서도 읽혀야 한다.
\definecolor{ink}{HTML}{1A1A1A}
\definecolor{soft}{HTML}{666666}
\definecolor{accent}{HTML}{7C3AED}
\definecolor{warn}{HTML}{B91C1C}

% 계층 색 — Simulink 모델의 블록 색과 같은 규칙을 쓴다
\definecolor{guid}{HTML}{1D4ED8}    % 유도
\definecolor{ctrl}{HTML}{EA580C}    % 제어
\definecolor{thr}{HTML}{D97706}     % 추진기
\definecolor{plant}{HTML}{16A34A}   % 운동모델
\definecolor{nav}{HTML}{7C3AED}     % 항법 · 신호처리
\definecolor{pathc}{HTML}{0D9488}   % 계획 경로 · 항적
\definecolor{errc}{HTML}{C2410C}    % 오차
\definecolor{flow}{HTML}{0369A1}    % 조류 · 바람

% ---- 선과 화살촉 -----------------------------------------------------------
\tikzset{
  sig/.style   = {ink, line width=0.5pt, -{Stealth[length=4pt, width=3pt]}},
  wire/.style  = {ink, line width=0.5pt},
  fb/.style    = {ink, line width=0.5pt, densely dashed,
                  -{Stealth[length=4pt, width=3pt]}},
  asig/.style  = {accent, line width=0.5pt, densely dashed,
                  -{Stealth[length=4pt, width=3pt]}},
  % 블록 — 흰 바탕, 직각, 검은 테두리. 색은 테두리에만 준다
  blk/.style   = {draw=ink, line width=0.55pt, fill=white,
                  minimum width=18mm, minimum height=8.5mm, inner sep=2pt, font=\small},
  gblk/.style  = {blk, draw=guid},
  cblk/.style  = {blk, draw=ctrl},
  tblk/.style  = {blk, draw=thr},
  pblk/.style  = {blk, draw=plant},
  nblk/.style  = {blk, draw=nav},
  % 합산점 · 분기점
  sum/.style   = {draw=ink, line width=0.55pt, fill=white, circle,
                  inner sep=0pt, minimum size=4.6mm, font=\scriptsize},
  dot/.style   = {ink, fill=ink, circle, inner sep=0pt, minimum size=1.5mm},
  % 글자
  sname/.style = {font=\small, inner sep=1.5pt},
  note/.style  = {font=\scriptsize, soft, inner sep=1.5pt},
  wnote/.style = {font=\scriptsize, warn, inner sep=1.5pt},
  anote/.style = {font=\scriptsize, accent, inner sep=1.5pt},
  axis/.style  = {ink, line width=0.5pt},
  tick/.style  = {font=\scriptsize, soft},
  % 묶음 상자 — 서브시스템 표시
  grp/.style   = {draw=soft, line width=0.4pt, densely dotted, rounded corners=1mm,
                  inner sep=3mm},
}
   (assets/src/ 안에서)
% 빌드:     bash _tools/tikz2svg.sh assets/src/w02-node-topic.tex assets/w02-node-topic.svg
%
% 왜 TikZ 인가
%   손으로 SVG 좌표를 쓰면 화살촉·선 굵기·글자 크기가 그림마다 미묘하게 달라진다.
%   그것이 "자료가 어설퍼 보인다"의 정체다. 스타일을 한 번만 정하고 상대좌표로
%   배치하면 좌표를 고쳐도 그림이 무너지지 않는다.
%   수식은 본문과 같은 TeX 글꼴로 나온다.
%
% 한글
%   ko.TeX(DVI 경로)를 쓴다. XeLaTeX 은 TikZ 그래픽을 PDF special 로 내보내는데
%   dvisvgm 이 그것을 받으려면 Ghostscript < 10.01 이 필요하다. 이 PC 에는
%   10.07 이 깔려 있어 그 경로가 막힌다. latex -> DVI -> dvisvgm 이 유일하게
%   동작하는 조합이다.

\usepackage[hangul]{kotex}
\usepackage{tikz}
\usepackage{amsmath}
\usepackage{xcolor}
\usetikzlibrary{arrows.meta, positioning, calc, fit, backgrounds, decorations.pathmorphing}

% ---- 색 --------------------------------------------------------------------
% 강조는 하나만. 나머지는 검정. 흑백 인쇄에서도 읽혀야 한다.
\definecolor{ink}{HTML}{1A1A1A}
\definecolor{soft}{HTML}{666666}
\definecolor{accent}{HTML}{7C3AED}
\definecolor{warn}{HTML}{B91C1C}

% 계층 색 — Simulink 모델의 블록 색과 같은 규칙을 쓴다
\definecolor{guid}{HTML}{1D4ED8}    % 유도
\definecolor{ctrl}{HTML}{EA580C}    % 제어
\definecolor{thr}{HTML}{D97706}     % 추진기
\definecolor{plant}{HTML}{16A34A}   % 운동모델
\definecolor{nav}{HTML}{7C3AED}     % 항법 · 신호처리
\definecolor{pathc}{HTML}{0D9488}   % 계획 경로 · 항적
\definecolor{errc}{HTML}{C2410C}    % 오차
\definecolor{flow}{HTML}{0369A1}    % 조류 · 바람

% ---- 선과 화살촉 -----------------------------------------------------------
\tikzset{
  sig/.style   = {ink, line width=0.5pt, -{Stealth[length=4pt, width=3pt]}},
  wire/.style  = {ink, line width=0.5pt},
  fb/.style    = {ink, line width=0.5pt, densely dashed,
                  -{Stealth[length=4pt, width=3pt]}},
  asig/.style  = {accent, line width=0.5pt, densely dashed,
                  -{Stealth[length=4pt, width=3pt]}},
  % 블록 — 흰 바탕, 직각, 검은 테두리. 색은 테두리에만 준다
  blk/.style   = {draw=ink, line width=0.55pt, fill=white,
                  minimum width=18mm, minimum height=8.5mm, inner sep=2pt, font=\small},
  gblk/.style  = {blk, draw=guid},
  cblk/.style  = {blk, draw=ctrl},
  tblk/.style  = {blk, draw=thr},
  pblk/.style  = {blk, draw=plant},
  nblk/.style  = {blk, draw=nav},
  % 합산점 · 분기점
  sum/.style   = {draw=ink, line width=0.55pt, fill=white, circle,
                  inner sep=0pt, minimum size=4.6mm, font=\scriptsize},
  dot/.style   = {ink, fill=ink, circle, inner sep=0pt, minimum size=1.5mm},
  % 글자
  sname/.style = {font=\small, inner sep=1.5pt},
  note/.style  = {font=\scriptsize, soft, inner sep=1.5pt},
  wnote/.style = {font=\scriptsize, warn, inner sep=1.5pt},
  anote/.style = {font=\scriptsize, accent, inner sep=1.5pt},
  axis/.style  = {ink, line width=0.5pt},
  tick/.style  = {font=\scriptsize, soft},
  % 묶음 상자 — 서브시스템 표시
  grp/.style   = {draw=soft, line width=0.4pt, densely dotted, rounded corners=1mm,
                  inner sep=3mm},
}
   (assets/src/ 안에서)
% 빌드:     bash _tools/tikz2svg.sh assets/src/w02-node-topic.tex assets/w02-node-topic.svg
%
% 왜 TikZ 인가
%   손으로 SVG 좌표를 쓰면 화살촉·선 굵기·글자 크기가 그림마다 미묘하게 달라진다.
%   그것이 "자료가 어설퍼 보인다"의 정체다. 스타일을 한 번만 정하고 상대좌표로
%   배치하면 좌표를 고쳐도 그림이 무너지지 않는다.
%   수식은 본문과 같은 TeX 글꼴로 나온다.
%
% 한글
%   ko.TeX(DVI 경로)를 쓴다. XeLaTeX 은 TikZ 그래픽을 PDF special 로 내보내는데
%   dvisvgm 이 그것을 받으려면 Ghostscript < 10.01 이 필요하다. 이 PC 에는
%   10.07 이 깔려 있어 그 경로가 막힌다. latex -> DVI -> dvisvgm 이 유일하게
%   동작하는 조합이다.

\usepackage[hangul]{kotex}
\usepackage{tikz}
\usepackage{amsmath}
\usepackage{xcolor}
\usetikzlibrary{arrows.meta, positioning, calc, fit, backgrounds, decorations.pathmorphing}

% ---- 색 --------------------------------------------------------------------
% 강조는 하나만. 나머지는 검정. 흑백 인쇄에서도 읽혀야 한다.
\definecolor{ink}{HTML}{1A1A1A}
\definecolor{soft}{HTML}{666666}
\definecolor{accent}{HTML}{7C3AED}
\definecolor{warn}{HTML}{B91C1C}

% 계층 색 — Simulink 모델의 블록 색과 같은 규칙을 쓴다
\definecolor{guid}{HTML}{1D4ED8}    % 유도
\definecolor{ctrl}{HTML}{EA580C}    % 제어
\definecolor{thr}{HTML}{D97706}     % 추진기
\definecolor{plant}{HTML}{16A34A}   % 운동모델
\definecolor{nav}{HTML}{7C3AED}     % 항법 · 신호처리
\definecolor{pathc}{HTML}{0D9488}   % 계획 경로 · 항적
\definecolor{errc}{HTML}{C2410C}    % 오차
\definecolor{flow}{HTML}{0369A1}    % 조류 · 바람

% ---- 선과 화살촉 -----------------------------------------------------------
\tikzset{
  sig/.style   = {ink, line width=0.5pt, -{Stealth[length=4pt, width=3pt]}},
  wire/.style  = {ink, line width=0.5pt},
  fb/.style    = {ink, line width=0.5pt, densely dashed,
                  -{Stealth[length=4pt, width=3pt]}},
  asig/.style  = {accent, line width=0.5pt, densely dashed,
                  -{Stealth[length=4pt, width=3pt]}},
  % 블록 — 흰 바탕, 직각, 검은 테두리. 색은 테두리에만 준다
  blk/.style   = {draw=ink, line width=0.55pt, fill=white,
                  minimum width=18mm, minimum height=8.5mm, inner sep=2pt, font=\small},
  gblk/.style  = {blk, draw=guid},
  cblk/.style  = {blk, draw=ctrl},
  tblk/.style  = {blk, draw=thr},
  pblk/.style  = {blk, draw=plant},
  nblk/.style  = {blk, draw=nav},
  % 합산점 · 분기점
  sum/.style   = {draw=ink, line width=0.55pt, fill=white, circle,
                  inner sep=0pt, minimum size=4.6mm, font=\scriptsize},
  dot/.style   = {ink, fill=ink, circle, inner sep=0pt, minimum size=1.5mm},
  % 글자
  sname/.style = {font=\small, inner sep=1.5pt},
  note/.style  = {font=\scriptsize, soft, inner sep=1.5pt},
  wnote/.style = {font=\scriptsize, warn, inner sep=1.5pt},
  anote/.style = {font=\scriptsize, accent, inner sep=1.5pt},
  axis/.style  = {ink, line width=0.5pt},
  tick/.style  = {font=\scriptsize, soft},
  % 묶음 상자 — 서브시스템 표시
  grp/.style   = {draw=soft, line width=0.4pt, densely dotted, rounded corners=1mm,
                  inner sep=3mm},
}


\begin{document}
\begin{tikzpicture}[x=1mm, y=1mm]

% ===========================================================================
% (가) 추진기 하나 — 각도 한계는 나비넥타이 꼴이다
% ===========================================================================
\begin{scope}[shift={(0,0)}]

  \node[font=\small\bfseries, anchor=north west] at (-34,46)
       {(가) 추진기 하나의 각도 한계};

  % 낼 수 있는 영역 — +-45 도 두 쪽
  \fill[guid, opacity=0.12] (0,0) -- (45:30) arc (45:-45:30) -- cycle;
  \fill[guid, opacity=0.12] (0,0) -- (135:30) arc (135:225:30) -- cycle;

  % 축
  \draw[axis] (-33,0) -- (33,0);
  \draw[axis] (0,-33) -- (0,33);
  \node[font=\scriptsize, ink, anchor=west]  at (33.5,0)  {$F_x$};
  \node[font=\scriptsize, ink, anchor=south] at (0,33.5)  {$F_y$};

  % 경계선
  \draw[guid, line width=0.6pt] (-30,-30) -- (30,30);
  \draw[guid, line width=0.6pt] (-30,30) -- (30,-30);

  % 각도 표시
  \draw[ctrl, line width=0.5pt, -{Stealth[length=2.6pt,width=2pt]}]
       (14,0) arc (0:45:14);
  \node[font=\scriptsize, ctrl, anchor=west] at (13,8) {$\delta_{\max}=45^\circ$};

  % 예시 벡터
  \draw[thr, line width=1.2pt, -{Stealth[length=4pt,width=3pt]}] (0,0) -- (25:26);
  \node[font=\scriptsize, thr, anchor=south west] at (25:27) {$(F_x, F_y)$};

  \node[font=\scriptsize, guid, anchor=north] at (0,-36) {낼 수 있는 방향};
  \node[font=\scriptsize, soft, anchor=north] at (0,-41) {(흰 쪽은 서보가 못 돌아간다)};

  \node[anchor=north, font=\scriptsize, text width=70mm, align=center] at (0,-48)
       {$\lvert \delta \rvert \le \delta_{\max}$
        \quad$\Longleftrightarrow$\quad
        $\lvert F_y \rvert \le \tan\delta_{\max}\,\lvert F_x \rvert$ \\[4pt]
        추력을 뒤집을 수 있으므로($T<0$ 허용) 왼쪽 영역도 함께 쓴다};
\end{scope}

% ===========================================================================
% (나) 요구 공간 — 쐐기 네 개
% ===========================================================================
\begin{scope}[shift={(96,0)}]

  \node[font=\small\bfseries, anchor=north west] at (-40,46)
       {(나) $N=0$ 단면 — pseudo-inverse 가 못 내는 방향};

  % 못 내는 쐐기 네 개
  \fill[warn, opacity=0.18] (0,0) -- (16.80:32) arc (16.80:37.33:32) -- cycle;
  \fill[warn, opacity=0.18] (0,0) -- (142.67:32) arc (142.67:163.20:32) -- cycle;
  \fill[warn, opacity=0.18] (0,0) -- (196.80:32) arc (196.80:217.33:32) -- cycle;
  \fill[warn, opacity=0.18] (0,0) -- (322.67:32) arc (322.67:343.20:32) -- cycle;

  \draw[warn, line width=0.5pt] (0,0) -- (16.80:32);
  \draw[warn, line width=0.5pt] (0,0) -- (37.33:32);
  \draw[warn, line width=0.5pt] (0,0) -- (142.67:32);
  \draw[warn, line width=0.5pt] (0,0) -- (163.20:32);
  \draw[warn, line width=0.5pt] (0,0) -- (196.80:32);
  \draw[warn, line width=0.5pt] (0,0) -- (217.33:32);
  \draw[warn, line width=0.5pt] (0,0) -- (322.67:32);
  \draw[warn, line width=0.5pt] (0,0) -- (343.20:32);

  \draw[soft, line width=0.3pt, densely dotted] (0,0) circle (32);

  % 축
  \draw[axis] (-36,0) -- (36,0);
  \draw[axis] (0,-36) -- (0,36);
  \node[font=\scriptsize, ink, anchor=west]  at (36.5,0) {$X$ (서지 힘)};
  \node[font=\scriptsize, ink, anchor=south] at (0,36.5) {$Y$ (스웨이 힘)};

  % 가능한 두 대표 방향
  \draw[plant, line width=1.2pt, -{Stealth[length=4pt,width=3pt]}] (0,0) -- (30,0);
  \node[font=\scriptsize, plant, anchor=north west] at (16,-1.5) {순수 전진};
  \draw[plant, line width=1.2pt, -{Stealth[length=4pt,width=3pt]}] (0,0) -- (0,30);
  \node[font=\scriptsize, plant, anchor=south east] at (-1.5,16)
       {순수 횡이동 ($\delta=\pm23.4^\circ$)};

  % 쐐기 각도 라벨
  \node[font=\scriptsize, warn, anchor=west] at (27:36) {$16.80^\circ \sim 37.33^\circ$};
  \node[font=\scriptsize, warn, anchor=east] at (207:35) {못 내는 쐐기};

  \node[anchor=north, font=\scriptsize, text width=72mm, align=center] at (0,-44)
       {요구한 방향이 쐐기 안에 들면 우현 추진기가 $45^\circ$ 를 넘겨야 한다 \\[3pt]
        \textbf{요구 공간이다 — 배를 위에서 본 그림이 아니다}};
\end{scope}

% ===========================================================================
% 아래 띠 — 왜 쐐기가 생기고, 어떻게 되살리는가
% ===========================================================================
\draw[wire, soft] (-34,-64) -- (134,-64);

\node[anchor=north west, font=\scriptsize, text width=80mm] at (-34,-69)
     {\textbf{쐐기는 계산으로 나온다 — 그려 본 것이 아니다} \\[4pt]
      $N=0$ 일 때 pseudo-inverse 해는 \\[3pt]
      \hspace*{4mm}$F_{x,L} = \tfrac{1}{2}\!\left(X + \tfrac{L}{b}Y\right)$,
      \quad $F_{y,L} = \tfrac{Y}{2}$ \\[2pt]
      \hspace*{4mm}$F_{x,R} = \tfrac{1}{2}\!\left(X - \tfrac{L}{b}Y\right)$,
      \quad $F_{y,R} = \tfrac{Y}{2}$ \\[4pt]
      우현에 $\lvert F_{y,R}\rvert \le \lvert F_{x,R}\rvert$ 를 넣고 풀면 \\[3pt]
      \hspace*{4mm}$1.3111 < X/Y < 3.3111$ \; 에서 \textbf{깨진다} \\[3pt]
      $L/b = 2.3110$ 이라는 기하 하나가 이 두 수를 정한다};

\node[anchor=north west, font=\scriptsize, text width=80mm] at (54,-69)
     {\textbf{null space 보정이 상당수를 되살린다} \\[4pt]
      같은 $\boldsymbol\tau$ 를 내는 해가 직선 하나를 이룬다 \\[3pt]
      \hspace*{4mm}$\mathbf{u}_e(\lambda) = \mathbf{B}^{+}\boldsymbol\tau
        + \lambda\,\mathbf{n}$,
      \quad $\mathbf{n} = [\,0,\ 1,\ 0,\ -1\,]^{\top}$ \\[4pt]
      $\lambda$ 는 좌우의 횡력만 주고받게 한다 — $X, Y, N$ 은 한 톨도 변하지 않는다.
      한계가 $\lambda$ 에 대해 1차식이라 허용 구간
      $[\lambda_{\mathrm{lo}},\ \lambda_{\mathrm{hi}}]$ 가 닫힌식으로 나온다 \\[4pt]
      \textbf{격자 실측} — 낼 수 있는 비율 $19.7\,\% \;\to\; 25.3\,\%$};

\node[anchor=north west, font=\scriptsize, warn, text width=168mm] at (-34,-114)
     {각도 한계를 $\pm15^\circ$ 로 좁히면 쐐기가 훨씬 넓어진다 — 기준 환경 실측에서
      각도 포화가 시간의 $80.7\,\%$, 못 낸 힘이 $309.97$ N 까지 늘었다.
      $\pm45^\circ$ 는 Sarda 등(2016)이 실제 WAM-V 에 달았던 서보의 한계값이다};

\end{tikzpicture}
\end{document}
